% Lab 5: BLoC
% Application Development for Mobile Devices
\documentclass[10pt,aspectratio=169,t]{beamer}
\usetheme[lab]{upbminimal}

\course{Application Development for Mobile Devices}
\decklabel{Lab 5}
\title{BLoC}
\subtitle{Cubit, events, and one test that proves the state}
\author{Dinu-Ștefan Rusu}
\institute{FILS, Universitatea Politehnica București}
\date{Week 10}

\begin{document}

\titleframe

\begin{frame}[eyebrow=Today]{What this lab is for}
  \vskip 2mm
  \begin{grid}[2]
    \cell[\faLayerGroup]{Separate logic from UI}{Move the todo list out of setState and into a Cubit, then a Bloc, while the screen only reacts to state.}
    \cell[\faVial]{Prove it with a test}{Write one bloc\_test that drives events through the bloc and checks the exact states it emits, in order.}
  \end{grid}
  \vskip 6mm
  \taskmeta{Time}{2 hours}{Submit}{Push to your admd-labs repo, branch \code{lab-5}, before next lab}
\end{frame}

\begin{frame}[fragile,eyebrow=Setup]{Where your code goes}
  \begin{steps}
    \item Open the project from Lab 4 and create a branch: \code{git checkout -b lab-5}
    \item Keep \code{lib/models/todo.dart} from Lab 4. Create a new \code{lib/todos/} folder for everything else.
    \item Run the app once before changing anything. If it does not build, fix that first.
  \end{steps}
  \begin{hint}
    Hot restart after wrapping the app in \code{BlocProvider}: a hot reload can leave a stale provider in the tree.
  \end{hint}
\end{frame}

\begin{frame}[fragile,eyebrow=Setup]{The feature's folder structure}
  \begin{codeblock}
lib/
  models/todo.dart
  todos/
    todo_repository.dart
    cubit/todo_cubit.dart
    cubit/todo_state.dart
    bloc/todo_bloc.dart
    bloc/todo_event.dart
    bloc/todo_state.dart
  screens/todo_screen.dart
  main.dart
test/todos/todo_bloc_test.dart
  \end{codeblock}
\end{frame}

\begin{frame}[fragile,eyebrow=Setup]{Packages to add}
\begin{shell}
flutter pub add flutter_bloc equatable
flutter pub add dev:bloc_test dev:mocktail
\end{shell}
  \begin{steps}
    \item \code{flutter\_bloc}: Cubit, Bloc, and the widgets that read them.
    \item \code{equatable}: value equality for state classes.
    \item \code{bloc\_test}: drives a bloc and asserts what it emits.
    \item \code{mocktail}: fakes the repository in the test.
  \end{steps}
  \begin{hint}
    \code{flutter pub add} writes the resolved version into \code{pubspec.yaml}. Commit that change together with your code.
  \end{hint}
\end{frame}

\begin{frame}[eyebrow=Task I]{TodoCubit: state without events}
  \begin{columns}[T]
    \begin{column}{0.55\textwidth}
      \begin{steps}
        \item Create \code{TodoState extends Equatable} holding the todo list and a loading flag.
        \item Create \code{TodoCubit extends Cubit<TodoState>} with \code{load}, \code{add}, \code{toggle}, \code{remove}.
        \item Wrap the app in \code{BlocProvider} whose \code{create} returns \code{TodoCubit()..load()}.
        \item Replace every \code{setState} in \code{TodoScreen} with a \code{BlocBuilder}.
      \end{steps}
    \end{column}
    \begin{column}{0.41\textwidth}
      \kv{Done when}{}
      \begin{checklist}
        \item No \code{setState} remains in \code{TodoScreen}.
        \item Toggling calls \code{TodoCubit}'s \code{toggle(id)} via \code{context.read}.
        \item Two states with the same todos compare equal.
      \end{checklist}
    \end{column}
  \end{columns}
  \begin{takeaway}{A cubit method and a bloc event both end in emit.}
    Task II only changes how the call arrives, not what happens after it.
  \end{takeaway}
\end{frame}

\begin{frame}[fragile,eyebrow=Task I]{The cubit's state}
\begin{dart}
class TodoState extends Equatable {
  final List<Todo> todos;
  final bool loading;
  const TodoState({this.todos = const [], this.loading = false});
  TodoState copyWith({List<Todo>? todos, bool? loading}) => TodoState(
        todos: todos ?? this.todos,
        loading: loading ?? this.loading,
      );
  @override
  List<Object?> get props => [todos, loading];
}
\end{dart}
\end{frame}

\begin{frame}[fragile,eyebrow=Task I]{TodoCubit: four methods, one shape}
\begin{dart}
class TodoCubit extends Cubit<TodoState> {
  TodoCubit() : super(const TodoState());
  Future<void> load() async {
    emit(state.copyWith(loading: true));
    emit(TodoState(todos: await fetchTodos()));
  }
  void toggle(int id) => emit(state.copyWith(todos: [
        for (final t in state.todos)
          if (t.id == id) t.copyWith(completed: !t.completed) else t,
      ]));
}
\end{dart}
\end{frame}

\begin{frame}[fragile,eyebrow=Task I]{Wiring: BlocProvider and BlocBuilder}
\begin{dart}
void main() => runApp(BlocProvider(
      create: (_) => TodoCubit()..load(),
      child: const TodoApp(),
    ));
// inside TodoScreen.build
BlocBuilder<TodoCubit, TodoState>(
  builder: (context, state) => state.loading
      ? const CircularProgressIndicator()
      : TodoList(
          todos: state.todos,
          onToggle: context.read<TodoCubit>().toggle,
        ),
)
\end{dart}
\end{frame}

\begin{frame}[eyebrow=Reference]{From Cubit to Bloc: why promote}
  \vskip 2mm
  \begin{center}
  \begin{tabular}{lll}
    \toprule
    & \thead{Cubit} & \thead{Bloc} \\
    \midrule
    Trigger & a direct method call & an event, via \code{add()} \\
    Multiple sources & harder to compose & one \code{on<Event>} per source \\
    Test grain & one call per case & one event per case \\
    Talks to the API & directly, or not at all & through a \code{TodoRepository} \\
    This lab & Task I & Task II onward \\
    \bottomrule
  \end{tabular}
  \end{center}
  \begin{takeaway}{Promote when one state can change for more than one reason.}
    Naming each reason as an event keeps the bloc's history readable and testable.
  \end{takeaway}
\end{frame}

\begin{frame}[eyebrow=Task II]{TodoBloc: events and a repository}
  \begin{columns}[T]
    \begin{column}{0.55\textwidth}
      \begin{steps}
        \item Define \code{TodoEvent} subclasses: \code{LoadTodos}, \code{AddTodo}, \code{ToggleTodo}, \code{RemoveTodo}.
        \item Wrap Lab 4's \code{fetchTodos()} inside a \code{TodoRepository} class.
        \item Create \code{TodoBloc}, a \code{Bloc<TodoEvent, TodoState>}, with one \code{on<Event>} handler per event.
        \item Swap \code{TodoCubit} for \code{TodoBloc} in \code{BlocProvider} and dispatch with \code{.add(...)}.
      \end{steps}
    \end{column}
    \begin{column}{0.41\textwidth}
      \kv{Done when}{}
      \begin{checklist}
        \item Every action calls \code{.add(...)}, never a cubit method.
        \item Only \code{TodoRepository} calls \code{dio}.
        \item Toggling and removing update the visible list.
      \end{checklist}
    \end{column}
  \end{columns}
  \begin{takeaway}{The repository is the only class that talks to the network.}
    TodoBloc receives todos; it never builds a Dio client itself.
  \end{takeaway}
\end{frame}

\begin{frame}[fragile,eyebrow=Task II]{The event hierarchy}
\begin{dart}
sealed class TodoEvent {}

class LoadTodos extends TodoEvent {}

class AddTodo extends TodoEvent {
  AddTodo(this.title);
  final String title;
}

// ToggleTodo(id) and RemoveTodo(id) carry the
// same single int id field as AddTodo carries title.
\end{dart}
\end{frame}

\begin{frame}[fragile,eyebrow=Task II]{TodoRepository: the only class that calls dio}
\begin{dart}
class TodoRepository {
  TodoRepository(this._dio);
  final Dio _dio;
  static const _url = 'https://jsonplaceholder.typicode.com/todos';

  Future<List<Todo>> fetchTodos() async {
    final res = await _dio.get(_url);
    return (res.data as List)
        .map((json) => Todo.fromJson(json))
        .toList();
  }
}
\end{dart}
\end{frame}

\begin{frame}[fragile,eyebrow=Task II]{TodoBloc: one handler per event}
\begin{dart}
class TodoBloc extends Bloc<TodoEvent, TodoState> {
  TodoBloc(this._repo) : super(const TodoState()) {
    on<LoadTodos>(_onLoad);
  }
  final TodoRepository _repo;
  Future<void> _onLoad(LoadTodos e, Emitter<TodoState> emit) async {
    emit(TodoState(todos: await _repo.fetchTodos()));
  }
}
\end{dart}
  \begin{hint}
    \code{on<ToggleTodo>}, \code{on<AddTodo>} and \code{on<RemoveTodo>} register the same way.
  \end{hint}
\end{frame}

\begin{frame}[eyebrow=Task III]{Proving it with bloc\_test}
  \begin{columns}[T]
    \begin{column}{0.55\textwidth}
      \begin{steps}
        \item Add \code{test/todos/todo\_bloc\_test.dart}.
        \item Fake the repository: extend \code{Mock} and implement \code{TodoRepository}.
        \item Use \code{blocTest} to add a todo, then toggle it, and list both states in \code{expect}.
        \item Run \code{flutter test} and confirm it is green.
      \end{steps}
    \end{column}
    \begin{column}{0.41\textwidth}
      \kv{Done when}{}
      \begin{checklist}
        \item The test builds \code{TodoBloc} with the fake, never the real repository.
        \item \code{act} fires \code{AddTodo} then \code{ToggleTodo}, in that order.
        \item \code{expect} lists both intermediate states, not just the last.
      \end{checklist}
    \end{column}
  \end{columns}
\end{frame}

\begin{frame}[fragile,eyebrow=Task III]{One test, two emitted states}
\begin{dart}
void main() {
  final repo = FakeTodoRepository();
  blocTest<TodoBloc, TodoState>(
    'add then toggle emits two states in order',
    build: () => TodoBloc(repo),
    act: (bloc) => bloc
      ..add(AddTodo('Milk'))
      ..add(ToggleTodo(1)),
    expect: () => [
      TodoState(todos: [Todo(id: 1, title: 'Milk')]),
      TodoState(todos: [Todo(id: 1, title: 'Milk', completed: true)]),
    ],
  );
}
\end{dart}
\end{frame}

\begin{frame}[eyebrow=Troubleshooting]{Errors you will see, and the fix}
  \vskip 2mm
  \trouble{Bad state: Cannot emit new states after calling close}{An async call finished after the bloc closed. Guard it, or cancel it in \code{close()}.}
  \trouble{BlocProvider.of() called with a context that does not contain a TodoBloc}{The reading widget sits above the BlocProvider. Move BlocProvider up the tree.}
  \trouble{MissingStubError: 'fetchTodos'}{The mocktail fake was never stubbed with \code{when(...)} before the test called it.}
  \trouble{Toggling a todo does nothing on screen}{The new state is \code{==} to the old one. Check that \code{copyWith} swaps the toggled item.}
\end{frame}

\closingframe{Push before you leave.}{Branch lab-5 · TodoBloc replaces TodoCubit · one passing bloc\_test in the PR}

\end{document}
